\section{Statement of the theorem}\label{sec:statement}

Let $C,\eta>0$ be fixed and let $T$ be sufficiently large in terms of these
parameters.  Put
\[
                         \Delta=\frac{T}{\log T}.
\]
Throughout, $w\in C_c^\infty(\R)$ is supported in $[T/4,2T]$ and satisfies
\begin{equation}\label{eq:w-derivatives}
                  w^{(j)}(t)\ll_j\Delta^{-j}
                  \qquad (j\geq0).
\end{equation}
The implied constants below may depend on $\eta,C,\eps$, on the constants in
the coefficient bounds, and on finitely many seminorms in
\eqref{eq:w-derivatives}.

For positive integers $h,k$ and complex $A,B,C_1$, define first in a region
of absolute convergence
\begin{equation}\label{eq:Z-series}
 Z_{A,B,C_1}(h,k)
 =\sum_{kab=hc}a^{-1/2-A}b^{-1/2-B}c^{-1/2-C_1}.
\end{equation}
The letters $a,b,c$ in this sum are local summation variables and are unrelated
to the coefficient sequences in the theorem.  To make the continuation
explicit, write $g=(h,k)$, $h=gh_0$, $k=gk_0$, and set
\[
                    q_Z=1+A+C_1,\qquad v_Z=1+B+C_1.
\]
Then
\begin{align}
 Z_{A,B,C_1}(h,k)
 &=h_0^{1/2+C_1}k_0^{-1/2-C_1}
   \zeta(q_Z)\zeta(v_Z) \notag\\
 &\quad\times
 \prod_{p^\nu\parallel h_0}
 (1-p^{-q_Z})(1-p^{-v_Z})
 \sum_{i+j\geq\nu}p^{-q_Zi-v_Zj}.
 \label{eq:Z-euler}
\end{align}
We use the right side as the meromorphic continuation of
\eqref{eq:Z-series}.  Its only shift poles are the two displayed zeta poles.

The exact archimedean factor is
\begin{equation}\label{eq:X-def}
 X_{A,C_1}(t)=\pi^{A+C_1}
 \frac{
 \Gamma\bigl((\frac12-A-it)/2\bigr)
 \Gamma\bigl((\frac12-C_1+it)/2\bigr)}{
 \Gamma\bigl((\frac12+A+it)/2\bigr)
 \Gamma\bigl((\frac12+C_1-it)/2\bigr)}.
\end{equation}

For two independent complex sequences $(a_h)$ and $(b_k)$, put
\begin{align}
 \I_{\alpha,\beta,\gamma}(a,b)
 ={}&\sum_{h\leq H}\sum_{k\leq K}
 \frac{a_h\overline{b_k}}{\sqrt{hk}}
 \int_{\R}w(t)
 \zeta\!\left(\tfrac12+\alpha+it\right)
 \zeta\!\left(\tfrac12+\beta+it\right) \notag\\
 &\hspace{34mm}\times
 \zeta\!\left(\tfrac12+\gamma-it\right)
 \left(\frac hk\right)^{it}\dd t.
 \label{eq:I-def}
\end{align}

\begin{theorem}[Fixed-Weil bilinear twisted cubic moment]
\label{thm:main}
Suppose $H,K\geq1$ and
\begin{equation}\label{eq:main-range}
                         H^{3/2}K\leq T^{1-\eta}.
\end{equation}
Let
\begin{equation}\label{eq:coeff-bounds}
             |a_h|\ll_\eps h^\eps,\qquad
             |b_k|\ll_\eps k^\eps
\end{equation}
for every $\eps>0$.  Uniformly for complex shifts satisfying
\begin{equation}\label{eq:shift-range}
              |\alpha|+|\beta|+|\gamma|\leq \frac{C}{\log T},
\end{equation}
one has
\begin{align}
 \I_{\alpha,\beta,\gamma}(a,b)
 ={}&\sum_{h\leq H}\sum_{k\leq K}
 \frac{a_h\overline{b_k}}{\sqrt{hk}}
 \int_{\R}w(t)\bigl\{
 Z_{\alpha,\beta,\gamma}(h,k) \notag\\
 &\quad+X_{\alpha,\gamma}(t)
 Z_{-\gamma,\beta,-\alpha}(h,k)
 +X_{\beta,\gamma}(t)
 Z_{\alpha,-\gamma,-\beta}(h,k)
 \bigr\}\dd t \notag\\
 &+O_{\eta,C,\eps}\!\left(
 T^\eps\{KH^{3/2}+T^{1/2}(HK)^{1/2}\}\right).
 \label{eq:main-asymptotic}
\end{align}
The expression in braces is its jointly holomorphic continuation across
\[
 \alpha+\gamma=0,\qquad \beta+\gamma=0,\qquad \alpha-\beta=0,
\]
including their common collision locus $\alpha=\beta=-\gamma$.
\end{theorem}

The theorem is formulated with exact gamma ratios.  Uniform Stirling
asymptotics give
\begin{equation}\label{eq:X-stirling}
       X_{A,C_1}(t)=\left(\frac{t}{2\pi}\right)^{-A-C_1}
                    \{1+O_C(T^{-1})\}
\end{equation}
on the support of $w$.  Consequently the two exact factors in
\eqref{eq:main-asymptotic} may be replaced by the corresponding powers used
in Bui's formulation.  The total replacement error is
$O(T^\eps(HK)^{1/2})$.

\begin{corollary}\label{cor:exponent-region}
If $H\leq T^{\theta_1}$ and $K\leq T^{\theta_3}$ range over a compact subset
of
\[
                         \frac32\theta_1+\theta_3<1,
\]
then \eqref{eq:main-asymptotic} has a power-saving error.  More precisely,
under \eqref{eq:main-range}, choosing the harmless $T^\eps$ loss with
$\eps<\eta/4$ gives an error $O(T^{1-\eta/4})$.
\end{corollary}

\begin{remark}\label{rem:independent}
No relation between $(a_h)$ and $(b_k)$ is used.  Complex conjugation on
$b_k$ is only notational; analytically the theorem is bilinear in two
independent sequences.
\end{remark}
